% Motor: XeLaTeX, fijado en .latexmkrc y en la receta del editor.
\documentclass[aspectratio=169,10pt,c]{beamer}

% !TEX root = slides.tex
% Preambulo comun: estilo UTP, colores, fuentes, plantillas de beamer y
% comandos de apoyo didactico. Se carga desde slides.tex tras \documentclass.

% La opción "c" de beamer reparte el espacio libre 1:1.5 (arriba:abajo), por lo
% que el contenido queda por encima del centro. Igualamos ambos rellenos.
\makeatletter
\define@key{beamerframe}{c}[true]{%
  \beamer@frametopskip=0pt plus 1fill\relax%
  \beamer@framebottomskip=0pt plus 1fill\relax%
  \beamer@frametopskipautobreak=0pt plus .5\paperheight\relax%
  \beamer@framebottomskipautobreak=0pt plus .5\paperheight\relax%
}
\makeatother

\usepackage{graphicx}
\usepackage{tikz}
\usepackage{xcolor}
\usepackage{iftex}

\ifPDFTeX
  \usepackage[utf8]{inputenc}
  \usepackage[T1]{fontenc}
  \usepackage{lmodern}          % tipografia escalable para el respaldo con pdflatex
  \renewcommand{\familydefault}{\sfdefault}
\else
  \usepackage{fontspec}
  \IfFontExistsTF{Arial}
    {\setsansfont{Arial}[BoldFont={Arial Bold}]}
    {\setsansfont{Helvetica}[BoldFont={Helvetica Bold}]}
\fi

\usepackage[spanish,es-nodecimaldot]{babel}
\usepackage{microtype}
\usepackage{booktabs}
\usepackage{amssymb}
\usepackage{array}
\usetikzlibrary{arrows.meta,positioning,fit,calc,shapes.geometric}

\graphicspath{{assets/}}

% ---------- Identidad visual UTP ----------
\definecolor{utpRojo}{HTML}{C42B35}
\definecolor{utpNegro}{HTML}{111111}
\definecolor{utpGris}{HTML}{5F5F5F}
\definecolor{utpGrisMedio}{HTML}{9A9A9A}
\definecolor{utpGrisClaro}{HTML}{F1F1F1}
\definecolor{utpNaranja}{HTML}{C1670B}

\usetheme{default}
\setbeamercolor{normal text}{fg=utpNegro,bg=white}
\setbeamercolor{structure}{fg=utpRojo}
\setbeamercolor{frametitle}{fg=utpNegro,bg=white}
\setbeamercolor{title}{fg=utpNegro}
\setbeamercolor{subtitle}{fg=utpNegro}
\setbeamercolor{author}{fg=utpNegro}
\setbeamercolor{date}{fg=utpNegro}
\setbeamercolor{block title}{fg=white,bg=utpNegro}
\setbeamercolor{block body}{fg=utpNegro,bg=utpGrisClaro}
\setbeamercolor{alertblock title}{fg=white,bg=utpRojo}
\setbeamercolor{alertblock body}{fg=utpNegro,bg=utpRojo!7}
\setbeamercolor{exampleblock title}{fg=white,bg=utpGris}
\setbeamercolor{exampleblock body}{fg=utpNegro,bg=utpGris!8}

\setbeamerfont{title}{series=\bfseries,size=\fontsize{19.5}{22.5}\selectfont}
\setbeamerfont{subtitle}{size=\fontsize{11.5}{13.5}\selectfont}
\setbeamerfont{date}{series=\bfseries,size=\fontsize{11.5}{13.5}\selectfont}
\setbeamerfont{author}{size=\fontsize{10.5}{12.5}\selectfont}
\setbeamerfont{frametitle}{series=\bfseries,size=\large}
\setbeamerfont{block title}{series=\bfseries,size=\small}
\setbeamerfont{block body}{size=\normalsize}

\setbeamertemplate{navigation symbols}{}
\setbeamertemplate{headline}{}
\setbeamertemplate{blocks}[default]
\setbeamertemplate{itemize item}{\color{utpRojo}\small\raise.15ex\hbox{$\blacktriangleright$}}
\setbeamertemplate{itemize subitem}{\color{utpGris}\scriptsize$\bullet$}
\setbeamertemplate{enumerate item}{\color{utpRojo}\insertenumlabel.}
\setbeamertemplate{background canvas}{\color{white}\rule{\paperwidth}{\paperheight}}

% Título de diapositiva: negrita negra a la izquierda, logo UTP a la derecha
\setbeamertemplate{frametitle}{%
  \begin{tikzpicture}[remember picture,overlay]
    \node[anchor=north east,inner sep=0pt]
      at ([xshift=-7mm,yshift=-5.5mm]current page.north east)
      {\includegraphics[width=17mm]{utp-logo.png}};
  \end{tikzpicture}%
  \vspace{1.2mm}%
  \begin{minipage}[t]{.80\textwidth}\raggedright\insertframetitle\end{minipage}\par
  \vspace{1.1mm}{\color{utpRojo}\rule{\textwidth}{1.1pt}}
}

% Solo numeración discreta; la plantilla UTP no lleva pie de página
\setbeamertemplate{footline}{%
  \hfill{\color{utpGrisMedio}\scriptsize\insertframenumber}\hspace{6mm}\vspace{2.5mm}%
}

% Portada UTP
\setbeamertemplate{title page}{%
  \begin{tikzpicture}[remember picture,overlay]
    \node[anchor=center,inner sep=0pt] at (current page.center) {%
      \includegraphics[width=\paperwidth,height=\paperheight]{utp-cover-bg.jpeg}%
    };
    \node[anchor=north,align=center,text=utpNegro,inner sep=0pt,
      text width=145mm,font=\usebeamerfont{title}]
      at ([yshift=-25.2mm]current page.north) {\inserttitle};
    \node[anchor=north,align=center,text=utpNegro,inner sep=0pt,
      text width=142mm,font=\usebeamerfont{subtitle}]
      at ([yshift=-38.2mm]current page.north) {\insertsubtitle};
    \node[anchor=north,align=center,text=utpNegro,inner sep=0pt,
      font=\usebeamerfont{date}]
      at ([yshift=-50.4mm]current page.north) {\insertdate};
    \node[anchor=north,align=center,text=utpNegro,inner sep=0pt,
      font=\usebeamerfont{author}]
      at ([yshift=-58.4mm]current page.north) {\insertauthor};
  \end{tikzpicture}%
}

% ---------- Comandos de apoyo didáctico ----------
\newcommand{\concepto}[2]{%
  \begin{block}{#1}\large #2\end{block}}
\newcommand{\pregunta}[1]{%
  \begin{alertblock}{Decide y explica}#1\end{alertblock}}
\newcommand{\analogia}[1]{%
  \begin{exampleblock}{En palabras simples}\small #1\end{exampleblock}}
\newcommand{\enlarealidad}[1]{%
  \begin{block}{Cómo se ve en la realidad}\small #1\end{block}}
\newcommand{\ojo}[1]{%
  \begin{alertblock}{Ojo con esto}\small #1\end{alertblock}}
\newcommand{\apoyo}{{\normalfont\footnotesize\color{utpGris}~· material de apoyo}}
\newcommand{\marca}[1]{\tikz[baseline=(m.base)]{%
  \node[rounded corners=1.5pt,fill=utpGrisClaro,draw=utpGrisMedio,inner xsep=3pt,inner ysep=1.5pt,
        font=\footnotesize] (m) {#1};}}

\tikzset{
  caja/.style={draw=utpNegro, very thick, rounded corners=1.5pt, fill=white,
    align=center, minimum height=0.72cm, inner sep=5pt, font=\small},
  cajaazul/.style={caja, draw=utpGris, fill=white},
  cajaverde/.style={caja, draw=utpNegro, fill=utpGrisClaro},
  cajanaranja/.style={caja, draw=utpRojo!75, fill=utpRojo!7},
  flecha/.style={-{Stealth[length=2.3mm]}, very thick, draw=utpNegro},
  flechavuelta/.style={{Stealth[length=2.3mm]}-{Stealth[length=2.3mm]}, very thick, draw=utpNegro}
}

% ---------- Notación pata de gallo (crow's foot) ----------
% La rotación indica hacia dónde está la entidad: 0 derecha, 180 izquierda,
% 90 arriba, -90 abajo. La pata se abre hacia la entidad del lado "muchos".
\tikzset{
  relacion/.style={thick, draw=utpNegro},
  pics/uno/.style={code={%
    \draw[thick,utpNegro,line cap=round] (0,-1.5mm) -- (0,1.5mm);}},
  pics/muchos/.style={code={%
    \draw[thick,utpNegro,line cap=round,line join=round]
      (2.6mm,1.5mm) -- (0,0) -- (2.6mm,-1.5mm);
    \draw[thick,utpNegro,line cap=round] (0,0) -- (2.6mm,0);}},
}
\newcommand{\simbuno}{\tikz[baseline=-0.6ex]{\draw[relacion](0,0)--(7mm,0);\pic at (4.4mm,0){uno};}}
\newcommand{\simbmuchos}{\tikz[baseline=-0.6ex]{\draw[relacion](0,0)--(7mm,0);\pic at (4.4mm,0){muchos};}}

\usepackage{colortbl}

\title{Subconsultas}
\subtitle{Consultas dentro de consultas,\\ con un repaso de servicios, clientes y conexión}
\date{Semana 12 - Sesión 23}
\author{Rel Guzman}

% Nombre de tabla, siempre en naranja
\newcommand{\tbl}[1]{\textcolor{utpNaranja}{\texttt{#1}}}
% Palabra reservada de SQL dentro de un bloque de código
\newcommand{\kw}[1]{\textcolor{utpRojo}{\textbf{#1}}}
% Comentario dentro de un bloque de código SQL
\newcommand{\com}[1]{\textcolor{utpGris}{-- #1}}
% Texto que escribe el usuario en la terminal
\newcommand{\tu}[1]{\textcolor{utpRojo}{\textbf{#1}}}
% Insignia de práctica
\newcommand{\ins}[1]{\tikz[baseline=-0.5ex]{\node[rounded corners=2pt,fill=utpRojo,text=white,
  font=\bfseries\scriptsize,inner xsep=5pt,inner ysep=2.5pt] {#1};}}
% Etiqueta oscura: nombre del gestor o del cliente
\newcommand{\etiqueta}[1]{\tikz[baseline=-0.5ex]{\node[rounded corners=2pt,fill=utpNegro,text=white,
  font=\bfseries\scriptsize,inner xsep=6pt,inner ysep=3pt] {#1};}}
% Número de paso
\newcommand{\npieza}[1]{\tikz[baseline=-0.6ex]{\node[circle,fill=utpRojo,text=white,
  font=\bfseries\scriptsize,inner sep=1.6pt,minimum size=4.2mm] {#1};}}
% Celda resaltada en rojo claro
\newcommand{\mal}[1]{{\setlength{\fboxsep}{1.2pt}\colorbox{utpRojo!18}{#1}}}
% La subconsulta, resaltada en amarillo dentro de un bloque de código
\newcommand{\dentro}[1]{{\setlength{\fboxsep}{1.5pt}\colorbox{yellow!45}{#1}}}
% Fila de una tabla que cumple la condición
\newcommand{\fila}{\rowcolor{yellow!45}}
% Marcas de sí y no al lado de una fila
\newcommand{\si}{$\checkmark$}
\newcommand{\no}{\textcolor{utpRojo}{$\times$}}
% Barra invertida de las rutas de Windows
\newcommand{\bs}{\textbackslash{}}
% Bloque de código sobre fondo gris
\newsavebox{\sqlbuf}
\newenvironment{sqlcode}[1][.86\textwidth]
  {\begin{center}\begin{lrbox}{\sqlbuf}\begin{minipage}{#1}\ttfamily\footnotesize\vspace{6pt}\raggedright}
  {\vspace{6pt}\end{minipage}\end{lrbox}\colorbox{utpGrisClaro}{\usebox{\sqlbuf}}\end{center}}
% Sangría dentro de un bloque de código
\newcommand{\ind}{\hspace*{1.6em}}
% Anatomía de una instrucción: cada pieza del código con su pregunta al lado
\tikzset{
  pieza/.style={fill=utpGrisClaro,font=\ttfamily\footnotesize,anchor=west,inner sep=5pt,
    minimum height=0.62cm},
  para/.style={font=\small,anchor=west,text=utpNegro}
}
\newcommand{\piezax}{7.6}
\newcommand{\pieza}[3]{%
  \node[pieza] at (0,#1) {#2};
  \node[para] at (\piezax,#1) {\textcolor{utpRojo}{$\blacktriangleleft$}\ #3};}
% Diapositiva de sección: una palabra bien grande (\seccion[0]{...} la deja sin línea roja)
\newcommand{\seccion}[2][1]{%
\begin{frame}[plain]
  \begin{tikzpicture}[remember picture,overlay]
    \node[anchor=north east,inner sep=0pt]
      at ([xshift=-10mm,yshift=-8mm]current page.north east)
      {\includegraphics[width=43mm]{utp-logo.png}};
    \node[anchor=center,inner sep=0pt] at (current page.center)
      {\resizebox{100mm}{!}{\bfseries\color{utpNegro}#2}};
    \ifnum#1=1
    \node[anchor=north,inner sep=0pt] at ([yshift=-11mm]current page.center)
      {\color{utpRojo}\rule{100mm}{1.6pt}};
    \fi
  \end{tikzpicture}
\end{frame}}
% Cabecera de columna dentro de una tabla en monoespaciado
\newcommand{\cab}[1]{\normalfont\bfseries #1}
% Resultado de una consulta: \res{columnas}{cabecera}{filas}
\newcommand{\res}[3]{%
  \centering{\scriptsize\color{utpGris}resultado}\\[3pt]
  \footnotesize\ttfamily
  \begin{tabular}{#1}
  \toprule
  #2 \\
  \midrule
  #3
  \bottomrule
  \end{tabular}}
% Tabla pequeña con su nombre arriba: \mini{nombre}{columnas}{cabecera}{filas}
\newcommand{\mini}[4]{%
  \begin{tabular}{@{}c@{}}
  \tbl{#1}\\[3pt]
  {\scriptsize\ttfamily
  \begin{tabular}{#2}
  \toprule
  #3 \\
  \midrule
  #4
  \bottomrule
  \end{tabular}}
  \end{tabular}}
% Forma de lo que devuelve una consulta: celdas con borde
\newcommand{\forma}[2]{{\ttfamily\footnotesize\renewcommand{\arraystretch}{1.25}%
  \begin{tabular}{#1}\hline #2\end{tabular}}}
% Tabla servicio de la peluquería, con dos filas que se pueden resaltar
\newcommand{\tablaservicio}[2]{%
  \mini{servicio}{rlr}{\cab{id} & \cab{nombre} & \cab{precio}}{%
  1 & Baño e higiene & 35.00 \\
  #1 2 & Corte de pelo & 55.00 \\
  3 & Corte de uñas & 15.00 \\
  #2 4 & Baño completo con corte & 80.00 \\}}
% Mensaje de ida (cliente a servidor) y de vuelta en un diagrama de secuencia
\newcommand{\ida}[2]{\draw[flecha] (0,#1)--node[above,font=\tiny,inner sep=1.5pt]{#2}(6.3,#1);}
\newcommand{\vuelta}[2]{\draw[flecha,draw=utpRojo] (6.3,#1)--node[above,font=\tiny,inner sep=1.5pt]{#2}(0,#1);}
% Captura de pantalla con números encima: \captura{ancho}{archivo}{marcadores}
% Dentro de los marcadores, (0,0) es la esquina inferior izquierda y (1,1) la superior derecha.
\newcommand{\captura}[3]{%
  \begin{tikzpicture}
    \node[anchor=south west,inner sep=0pt,draw=utpGrisMedio] (img) at (0,0) {\includegraphics[width=#1]{#2}};
    \begin{scope}[x={(img.south east)},y={(img.north west)}]
      \foreach \px/\py/\n in #3 \node[pin] at (\px,\py) {\n};
    \end{scope}
  \end{tikzpicture}}
% Explicación de un campo numerado de una captura
\newcommand{\explica}[3]{\par\noindent\npieza{#1}\ \textbf{#2}\par\vspace{2pt}{\leftskip=6.2mm #3\par}\vspace{10pt}}
% Estilos de los diagramas de esta sesión
\tikzset{
  num/.style={circle,fill=utpRojo,text=white,font=\bfseries\tiny,inner sep=1.2pt,minimum size=3.6mm},
  vida/.style={utpGrisMedio,thick,dashed},
  pin/.style={circle,fill=utpRojo,text=white,font=\bfseries\scriptsize,inner sep=1pt,
    minimum size=4.4mm,draw=white,thick},
  codigo/.style={font=\ttfamily\footnotesize,inner sep=3pt},
  sub/.style={codigo,draw=utpRojo,very thick,fill=yellow!45,rounded corners=2pt}
}
\begin{document}
{\setbeamertemplate{footline}{}\maketitle}
\begin{frame}[c]
\begin{block}{Logro de la sesión}
Al finalizar, podrás \textbf{conectarte} a un gestor sabiendo qué datos necesitas y escribir \textbf{subconsultas}: una consulta dentro de otra, para responder preguntas que se resuelven en \textbf{dos pasos}.
\end{block}
\end{frame}
% =======================================================================
% Revisión de la sesión anterior
% ---------------------------------------------------------------------
\begin{frame}{Revisión de la sesión anterior: el gestor es un servicio}
\begin{columns}[c,onlytextwidth]
\column{.47\textwidth}
\small El gestor es un programa que \textbf{no ves}. Funciona como \textbf{servicio}:
\vspace{0.2cm}
\begin{itemize}\setlength\itemsep{6pt}
  \item no tiene ventana,
  \item trabaja en \textbf{segundo plano},
  \item arranca solo al encender Windows,
  \item y se queda \textbf{esperando pedidos}.
\end{itemize}
\vspace{0.3cm}
\small Los ves con \textbf{Windows + R} $\rightarrow$ \texttt{services.msc}.
\column{.5\textwidth}
\centering
\begin{tikzpicture}[font=\scriptsize]
  \node[caja,text width=1.9cm] (wb) at (0,0) {MySQL\\Workbench};
  \node[caja,text width=1.9cm] (ss) at (2.6,0) {SSMS};
  \node[caja,cajanaranja,text width=1.9cm] (my) at (0,-2.3) {\texttt{\textbf{mysqld.exe}}\\{\tiny servicio \texttt{MySQL80}}};
  \node[caja,cajanaranja,text width=1.9cm] (sq) at (2.6,-2.3) {\texttt{\textbf{sqlservr.exe}}\\{\tiny servicio \texttt{SQL Server}}};
  \draw[flechavuelta] (wb)--(my); \draw[flechavuelta] (ss)--(sq);
  \draw[utpGrisMedio,thick,dashed] (-1.5,-1.15)--(4.1,-1.15);
  \node[anchor=south west,font=\scriptsize\bfseries] at (-1.5,0.5) {Clientes: piden};
  \node[anchor=north west,font=\scriptsize\bfseries,text=utpRojo] at (-1.5,-2.95) {Servicios: atienden};
\end{tikzpicture}
\end{columns}
\vspace{0.3cm}
\concepto{Servicio}{\small Un programa sin ventana que arranca con Windows y siempre está listo para atender. Si el servicio está apagado, \textbf{nadie} puede conectarse.}
\end{frame}
% ---------------------------------------------------------------------
\begin{frame}{Revisión: los clientes}
\concepto{Cliente}{\small \textbf{Cualquier programa} que se conecta al gestor para enviarle consultas y recibir resultados. Puede tener ventana (Workbench), ser una terminal (\texttt{mysql}) o no tener ventana (una app web, un programa en Node).}
\small Cada gestor tiene su cliente:
\begin{center}
\begin{tikzpicture}[font=\scriptsize]
  \foreach \i/\c/\g in {0/SSMS/SQL Server, 1/MySQL Workbench/MySQL,
      2/DB Browser for SQLite/SQLite, 3/pgAdmin/PostgreSQL} {
    \node[caja,text width=2.5cm] (c\i) at (3.3*\i,0) {\textbf{\c}};
    \node[caja,cajanaranja,text width=2.5cm] (g\i) at (3.3*\i,-1.2) {\g};
    \draw[flecha] (c\i)--(g\i);
  }
\end{tikzpicture}
\end{center}
\vspace{0.1cm}
\small Hay otros que sirven para \textbf{varios} gestores o para \textbf{dibujar} el modelo:\\[5pt]
\etiqueta{DataGrip} \etiqueta{DBeaver} \etiqueta{Luna Modeler} \etiqueta{Visual Paradigm} \etiqueta{HeidiSQL} {\footnotesize \ldots entre otros.}
\vspace{0.05cm}
\ojo{Si cierras el cliente, los datos \textbf{no} se borran: siguen en el gestor.}
\end{frame}
% ---------------------------------------------------------------------
\begin{frame}{Los cinco datos de una conexión}
\small ¿Qué necesitas para conectarte? Cinco datos:
\begin{center}
\renewcommand{\arraystretch}{1.4}
\begin{tabular}{@{}>{\bfseries}l>{\raggedright}p{5.0cm}>{\raggedright\arraybackslash}p{3.6cm}@{}}
\toprule
Dato & \bfseries Pregunta & \bfseries En clase \\
\midrule
Host & ¿en qué computadora? & \texttt{localhost} \\
Puerto & ¿a qué programa? & \texttt{3306} \\
Usuario & ¿quién eres? & \texttt{root} \\
Contraseña & ¿cómo lo demuestras? & la de la instalación \\
Base de datos & ¿con cuál empiezas? \textit{(opcional)} & \texttt{peluqueria} \\
\bottomrule
\end{tabular}
\end{center}
\vspace{0.3cm}
\concepto{Idea clave}{\small Todos los clientes piden \textbf{estos cinco datos}. Solo cambia dónde se escriben.}
\end{frame}
% ---------------------------------------------------------------------
\begin{frame}{Los cinco datos en SQL Server con Windows}
\begin{center}
\small
\renewcommand{\arraystretch}{1.35}
\begin{tabular}{@{}>{\bfseries}l>{\raggedright}p{3.0cm}c>{\raggedright\arraybackslash}p{7.2cm}@{}}
\toprule
Dato & \etiqueta{MySQL} & & \etiqueta{SQL Server con autenticación de Windows} \\
\midrule
Host & \texttt{localhost} & $\rightarrow$ & \textbf{Server Name}: \texttt{localhost} o \texttt{localhost\bs SQLEXPRESS} \\
Puerto & \texttt{3306} & $\rightarrow$ & no se escribe: usa \texttt{1433} solo \\
Usuario & \texttt{root} & $\rightarrow$ & \mal{tu usuario de Windows}: se llena solo \\
Contraseña & la de la instalación & $\rightarrow$ & \mal{no se pide} \\
Base de datos & \texttt{peluqueria} & $\rightarrow$ & \textbf{Database Name}, o luego \texttt{USE peluqueria;} \\
\bottomrule
\end{tabular}
\end{center}
\vspace{0.2cm}
\concepto{Autenticación de Windows}{\small Windows ya comprobó quién eres cuando iniciaste sesión en la laptop. SQL Server \textbf{confía} en eso: por eso no pide contraseña.}
\end{frame}
% ---------------------------------------------------------------------
\begin{frame}{Los cinco datos en SQLite}
\begin{center}
\small
\renewcommand{\arraystretch}{1.35}
\begin{tabular}{@{}>{\bfseries}l>{\raggedright}p{3.0cm}c>{\raggedright\arraybackslash}p{7.2cm}@{}}
\toprule
Dato & \etiqueta{MySQL} & & \etiqueta{SQLite} \\
\midrule
Host & \texttt{localhost} & $\rightarrow$ & \mal{no hay}: el archivo está en tu computadora \\
Puerto & \texttt{3306} & $\rightarrow$ & \mal{no hay}: no hay servidor esperando \\
Usuario & \texttt{root} & $\rightarrow$ & \mal{no hay} \\
Contraseña & la de la instalación & $\rightarrow$ & \mal{no hay} \\
Base de datos & \texttt{peluqueria} & $\rightarrow$ & \textbf{el archivo}: \texttt{peluqueria.db} \\
\bottomrule
\end{tabular}
\end{center}
\vspace{0.1cm}
\begin{columns}[T,onlytextwidth]
\column{.48\textwidth}
\centering\etiqueta{DB Browser for SQLite}\\[6pt]
\small \textit{Abrir base de datos} $\rightarrow$ eliges \texttt{peluqueria.db}
\column{.48\textwidth}
\centering\etiqueta{Terminal}
\begin{sqlcode}[.8\textwidth]
\tu{sqlite3 peluqueria.db}
\end{sqlcode}
\end{columns}
\vspace{0.1cm}
\concepto{Idea clave}{\small En SQLite, conectarse es \textbf{abrir un archivo}. Solo necesitas un dato: \textbf{dónde está}.}
\end{frame}
% ---------------------------------------------------------------------
\begin{frame}{¿Qué pasa al conectarse?}
\begin{center}
\begin{tikzpicture}[font=\scriptsize]
  \node[caja,text width=2cm] (c) at (0,0) {\textbf{Cliente}\\{\tiny Workbench}};
  \node[caja,cajanaranja,text width=2.3cm] (s) at (6.3,0) {\textbf{Servidor}\\{\tiny MySQL en localhost:3306}};
  \draw[vida] (c.south)--(0,-3.9); \draw[vida] (s.south)--(6.3,-3.9);
  \ida{-0.95}{1. llega a localhost, puerto 3306}
  \vuelta{-1.5}{2. «¿quién eres?»}
  \ida{-2.05}{3. usuario y contraseña}
  \vuelta{-2.6}{4. «aceptado»: conexión abierta}
  \ida{-3.15}{5. consultas \ldots}
  \vuelta{-3.7}{6. \ldots{} y resultados}
\end{tikzpicture}
\end{center}
\vspace{0.2cm}
\small Si falla el paso 1, el problema es el host o el puerto. Si falla el 4, el usuario o la contraseña.
\end{frame}
% ---------------------------------------------------------------------
\begin{frame}{Desde la terminal}
\small La \textbf{terminal} es una ventana para escribir órdenes (PowerShell o cmd).
\vspace{0.3cm}
\begin{center}
\renewcommand{\piezax}{3.5}
\begin{tikzpicture}
  \pieza{0}{mysql}{el programa cliente}
  \pieza{-0.68}{-h localhost}{host}
  \pieza{-1.36}{-P 3306}{puerto}
  \pieza{-2.04}{-u root}{usuario}
  \pieza{-2.72}{-p}{pide la contraseña}
  \pieza{-3.4}{peluqueria}{base de datos}
\end{tikzpicture}
\end{center}
\vspace{0.2cm}
\begin{sqlcode}[.7\textwidth]
\tu{mysql -h localhost -P 3306 -u root -p peluqueria}
\end{sqlcode}
\end{frame}
% ---------------------------------------------------------------------
\begin{frame}{Desde MySQL Workbench}
\begin{columns}[c,onlytextwidth]
\column{.58\textwidth}
\centering
\captura{\textwidth}{workbench-nueva-conexion.png}{{0.39/0.706/1, 0.61/0.706/2, 0.60/0.632/3}}
\column{.38\textwidth}
\explica{1}{Hostname}{el host: \texttt{127.0.0.1} es localhost}
\explica{2}{Port}{el puerto: \texttt{3306}}
\explica{3}{Username}{el usuario: \texttt{root}}
\vspace{0.2cm}
\small Luego \textbf{OK}. La contraseña se pide al conectar.
\end{columns}
\end{frame}
% ---------------------------------------------------------------------
\begin{frame}{Desde SSMS (SQL Server)}
\begin{columns}[c,onlytextwidth]
\column{.38\textwidth}
\centering
\captura{4.6cm}{ssms-connect.png}{{0.80/0.546/1, 0.80/0.506/2, 0.80/0.468/3}}
\column{.58\textwidth}
\explica{1}{Server Name}{el host: el nombre de tu computadora}
\explica{2}{Authentication}{\textit{Windows Authentication}: entras con tu cuenta de Windows}
\explica{3}{User Name}{el usuario; con Windows Authentication se llena solo}
\vspace{0.2cm}
\small Luego \textbf{Connect}. SQL Server no pide puerto: lo resuelve solo.
\end{columns}
\end{frame}
% ---------------------------------------------------------------------
\begin{frame}{Desde un programa}
\concepto{Cadena de conexión}{\small Los cinco datos escritos en \textbf{una sola línea}. Así se conectan las aplicaciones.}
\vspace{0.3cm}
\begin{center}
\begin{tikzpicture}[every node/.style={font=\ttfamily,inner sep=1pt}]
  \node (a) {mysql://};
  \node[anchor=west,text=utpRojo] (b) at (a.east) {root};
  \node[anchor=west] (c1) at (b.east) {:};
  \node[anchor=west,text=utpNaranja] (c) at (c1.east) {clave};
  \node[anchor=west] (d1) at (c.east) {@};
  \node[anchor=west,text=utpRojo] (d) at (d1.east) {localhost};
  \node[anchor=west] (e1) at (d.east) {:};
  \node[anchor=west,text=utpNaranja] (e) at (e1.east) {3306};
  \node[anchor=west] (f1) at (e.east) {/};
  \node[anchor=west,text=utpRojo] (f) at (f1.east) {peluqueria};
  \foreach \n/\t in {b/usuario,c/contraseña,d/host,e/puerto,f/base}
    \node[below=7pt of \n,font=\scriptsize\sffamily] {\t};
\end{tikzpicture}
\end{center}
\vspace{0.4cm}
\small Ejemplo en Node.js: \texttt{ejemplos/conexion\_al\_gestor/3\_cliente\_node/}.
\end{frame}
% ---------------------------------------------------------------------
\begin{frame}{Cuando no conecta}
\begin{center}
\footnotesize
\renewcommand{\arraystretch}{1.4}
\begin{tabular}{@{}>{\raggedright}p{6.4cm}>{\raggedright\arraybackslash}p{4.6cm}@{}}
\toprule
\bfseries Mensaje & \bfseries Qué revisar \\
\midrule
\texttt{ERROR 2003: Can't connect to MySQL server on 'localhost:3306'} & ¿el servicio está encendido? ¿host y puerto? \\
\texttt{ERROR 1045: Access denied for user 'root'@'localhost'} & usuario y contraseña \\
\texttt{ERROR 1049: Unknown database 'peluqueria'} & ¿existe la base? \\
\bottomrule
\end{tabular}
\end{center}
\vspace{0.3cm}
\concepto{¿Está encendido el servicio?}{\small \textbf{Windows + R} $\rightarrow$ \texttt{services.msc} $\rightarrow$ \texttt{MySQL80}. Debe decir \textbf{En ejecución}; si no, clic derecho $\rightarrow$ \textit{Iniciar}.}
\end{frame}
% =======================================================================
\seccion{Subconsultas}
\begin{frame}{Una pregunta que necesita otra pregunta}
\begin{center}
\begin{tikzpicture}
  \node[caja,cajanaranja,font=\large,inner sep=9pt] {¿Qué servicios cuestan \textbf{más que el promedio}?};
\end{tikzpicture}
\end{center}
\vspace{0.2cm}
\begin{columns}[c,onlytextwidth]
\column{.42\textwidth}
\centering
\tablaservicio{}{}
\column{.55\textwidth}
\small Para responderla, primero necesitas \textbf{otro dato}:
\vspace{0.3cm}
\begin{center}
\begin{tikzpicture}[font=\small]
  \node[caja,text width=2.3cm] (p1) at (0,0) {\npieza{1}\\¿cuál es el promedio?};
  \node[caja,text width=2.3cm] (p2) at (3.6,0) {\npieza{2}\\¿quiénes lo superan?};
  \draw[flecha] (p1)--node[above,font=\scriptsize\bfseries,text=utpRojo]{46.25}(p2);
\end{tikzpicture}
\end{center}
\vspace{0.2cm}
\small La respuesta de \npieza{1} se usa en \npieza{2}.
\end{columns}
\end{frame}
% ---------------------------------------------------------------------
\begin{frame}{Paso a paso: con dos consultas}
\begin{columns}[T,onlytextwidth]
\column{.49\textwidth}
\npieza{1}\ \small ¿Cuál es el promedio?
\begin{sqlcode}[\dimexpr\textwidth-2\fboxsep\relax]
\kw{SELECT AVG}(precio) \kw{FROM} servicio;
\end{sqlcode}
\vspace{-0.1cm}
\res{r}{\cab{AVG(precio)}}{46.25 \\}
\column{.49\textwidth}
\npieza{2}\ \small ¿Qué servicios cuestan más?
\begin{sqlcode}[\dimexpr\textwidth-2\fboxsep\relax]
\kw{SELECT} nombre, precio \kw{FROM} servicio\\
\kw{WHERE} precio > \mal{46.25};
\end{sqlcode}
\vspace{-0.1cm}
\res{lr}{\cab{nombre} & \cab{precio}}{Corte de pelo & 55.00 \\ Baño completo con corte & 80.00 \\}
\end{columns}
\vspace{0.3cm}
\ojo{Funciona, pero copiaste \texttt{46.25} \textbf{a mano}. Si mañana cambia un precio, ese número queda viejo.}
\end{frame}
% ---------------------------------------------------------------------
\begin{frame}{En una sola: la subconsulta}
\small En lugar de escribir el número, escribe \textbf{la consulta que lo calcula}:
\vspace{0.2cm}
\begin{center}
\begin{tikzpicture}
  \node[codigo,anchor=west] (a) at (0,0) {\kw{WHERE} precio >};
  \node[codigo,anchor=west,draw=utpRojo,thick,fill=utpRojo!18] (n) at ([xshift=2pt]a.east) {46.25};
  \node[codigo,anchor=west] (b) at (0,-1.5) {\kw{WHERE} precio >};
  \node[sub,anchor=west] (s) at ([xshift=2pt]b.east) {(\kw{SELECT AVG}(precio) \kw{FROM} servicio)};
  \draw[flecha] (n.south)--node[right,font=\scriptsize]{se reemplaza por}(n.south|-s.north);
\end{tikzpicture}
\end{center}
\vspace{0.1cm}
\begin{sqlcode}[.62\textwidth]
\kw{SELECT} nombre, precio\\
\kw{FROM} servicio\\
\kw{WHERE} precio > \dentro{(\kw{SELECT AVG}(precio) \kw{FROM} servicio)};
\end{sqlcode}
\vspace{-0.1cm}
\concepto{Subconsulta}{\small Una consulta escrita \textbf{dentro de otra}, entre \textbf{paréntesis}.}
\end{frame}
% ---------------------------------------------------------------------
\begin{frame}{Como en matemáticas: primero el paréntesis}
\begin{center}
\begin{tikzpicture}[font=\small]
  \tikzset{etapa/.style={caja,minimum height=0.9cm},
    paso/.style={flecha},
    lbl/.style={above,font=\scriptsize,text=utpGris,align=center}}
  % Matemáticas
  \node[font=\small\bfseries,anchor=east] at (-0.4,0) {Matemáticas};
  \node[etapa,text width=2.3cm] (m1) at (1.4,0) {$2 \times \colorbox{yellow!45}{$(3 + 4)$}$};
  \node[etapa,text width=1.8cm] (m2) at (5.7,0) {$2 \times 7$};
  \node[etapa,cajaverde,text width=2.1cm] (m3) at (10.0,0) {$14$};
  \draw[paso] (m1)--node[lbl]{primero el\\paréntesis}(m2);
  \draw[paso] (m2)--node[lbl]{luego lo\\de afuera}(m3);
  % SQL
  \node[font=\small\bfseries,anchor=east] at (-0.4,-1.8) {SQL};
  \node[etapa,text width=2.9cm,font=\ttfamily\scriptsize] (s1) at (1.4,-1.8)
    {precio >\\\dentro{(\kw{SELECT AVG}(precio)}\\\dentro{\kw{FROM} servicio)}};
  \node[etapa,text width=1.8cm,font=\ttfamily\scriptsize] (s2) at (5.7,-1.8) {precio > 46.25};
  \node[etapa,cajaverde,text width=2.1cm,font=\scriptsize] (s3) at (10.0,-1.8) {Corte de pelo\\Baño completo\ldots};
  \draw[paso] (s1)--node[lbl]{primero la\\de adentro}(s2);
  \draw[paso] (s2)--node[lbl]{luego la\\de afuera}(s3);
\end{tikzpicture}
\end{center}
\vspace{0.4cm}
\analogia{El gestor resuelve \textbf{primero} la consulta de adentro, \textbf{reemplaza} el paréntesis por su resultado y \textbf{después} ejecuta la de afuera.}
\end{frame}
% ---------------------------------------------------------------------
\begin{frame}{Una consulta dentro de otra}
\begin{center}
\begin{tikzpicture}
  \node[codigo,anchor=west] (o1) at (0,0) {\kw{SELECT} nombre, precio \kw{FROM} servicio};
  \node[codigo,anchor=west] (o2) at (0,-0.75) {\kw{WHERE} precio >};
  \node[sub,anchor=west,inner sep=6pt] (in) at ([xshift=3pt]o2.east) {(\kw{SELECT AVG}(precio) \kw{FROM} servicio)};
  \node[font=\small\bfseries,text=utpRojo,anchor=north] (lin) at ([yshift=-4pt]in.south) {subconsulta (la de adentro)};
  \node[draw=utpNegro,very thick,rounded corners=4pt,inner sep=12pt,fit=(o1)(o2)(in)(lin)] (out) {};
  \node[num,minimum size=5mm,font=\bfseries\scriptsize] at (in.north east) {1};
  \node[num,minimum size=5mm,font=\bfseries\scriptsize] at (out.north west) {2};
  \node[font=\small\bfseries,anchor=south west] at ([xshift=8pt,yshift=2pt]out.north west) {consulta principal (la de afuera)};
\end{tikzpicture}
\end{center}
\vspace{0.4cm}
\begin{itemize}\setlength\itemsep{8pt}
  \item \npieza{1} La \textbf{subconsulta} se resuelve primero y entrega un resultado: \texttt{46.25}.
  \item \npieza{2} La \textbf{consulta principal} usa ese resultado como si lo hubieras escrito tú.
\end{itemize}
\end{frame}
% ---------------------------------------------------------------------
\begin{frame}{¿Qué devuelve la de adentro?}
\small Lo primero es mirar la \textbf{forma} de su resultado. Eso decide cómo se usa.
\vspace{0.3cm}
\begin{columns}[T,onlytextwidth]
\column{.31\textwidth}
\centering
{\large\bfseries Un valor}\\[8pt]
\forma{|c|}{\cellcolor{yellow!45}46.25 \\ \hline}\\[8pt]
\footnotesize una fila, una columna\\[10pt]
\etiqueta{= \ > \ < }\\[6pt]
\scriptsize se compara
\column{.31\textwidth}
\centering
{\large\bfseries Una lista}\\[8pt]
\forma{|c|}{\cellcolor{yellow!45}1 \\ \hline \cellcolor{yellow!45}2 \\ \hline}\\[8pt]
\footnotesize una columna, varias filas\\[10pt]
\etiqueta{IN \ NOT IN}\\[6pt]
\scriptsize se busca dentro
\column{.31\textwidth}
\centering
{\large\bfseries Una tabla}\\[8pt]
\forma{|c|c|}{\cellcolor{yellow!45}1 & \cellcolor{yellow!45}3 \\ \hline \cellcolor{yellow!45}2 & \cellcolor{yellow!45}2 \\ \hline}\\[8pt]
\footnotesize varias columnas y filas\\[10pt]
\etiqueta{FROM}\\[6pt]
\scriptsize se usa como tabla
\end{columns}
\vspace{0.5cm}
\small Veamos un ejemplo de cada una.
\end{frame}
% ---------------------------------------------------------------------
\begin{frame}{Un valor: el servicio más caro}
\begin{columns}[c,onlytextwidth]
\column{.42\textwidth}
\centering
\tablaservicio{}{\fila}
\column{.55\textwidth}
\begin{sqlcode}[\dimexpr\textwidth-2\fboxsep\relax]
\kw{SELECT} nombre, precio\\
\kw{FROM} servicio\\
\kw{WHERE} precio = \dentro{(\kw{SELECT MAX}(precio) \kw{FROM} servicio)};
\end{sqlcode}
\vspace{0.1cm}
\begin{center}
\begin{tikzpicture}[font=\scriptsize]
  \node[caja,text width=2.4cm] (a) at (0,0) {\npieza{1} la de adentro\\devuelve \textbf{80.00}};
  \node[caja,text width=2.4cm] (b) at (3.4,0) {\npieza{2} la de afuera busca\\\texttt{precio = 80.00}};
  \draw[flecha] (a)--(b);
\end{tikzpicture}
\end{center}
\vspace{0.1cm}
\res{lr}{\cab{nombre} & \cab{precio}}{Baño completo con corte & 80.00 \\}
\end{columns}
\end{frame}
% ---------------------------------------------------------------------
\begin{frame}{Un valor de otra tabla: las mascotas de Lucía}
\small En \tbl{mascota} no está el nombre del dueño, solo su \texttt{cliente\_id}. \textbf{No sabes} el id de Lucía: que lo busque la subconsulta.
\vspace{0.2cm}
\begin{center}
\begin{tikzpicture}
  \node (cl) at (0,0) {\mini{cliente}{rl}{\cab{id} & \cab{nombre}}{%
    \fila 1 & Lucía Ramos \\
    2 & Pedro Salas \\
    3 & Rosa Medina \\}};
  \node (ma) at (7,0) {\mini{mascota}{llr}{\cab{nombre} & \cab{especie} & \cab{cliente\_id}}{%
    \fila Firulais & perro & 1 \\
    \fila Michi & gato & 1 \\
    Rocky & perro & 2 \\
    Luna & gato & 3 \\}};
  \draw[flecha] (cl.east)--node[above,font=\small\bfseries]{\npieza{1} id = 1}node[below,font=\scriptsize]{\npieza{2} busca el 1}(ma.west);
\end{tikzpicture}
\end{center}
\vspace{0.1cm}
\begin{sqlcode}[.84\textwidth]
\kw{SELECT} nombre, especie \kw{FROM} mascota\\
\kw{WHERE} cliente\_id = \dentro{(\kw{SELECT} id \kw{FROM} cliente \kw{WHERE} nombre = 'Lucía Ramos')};
\end{sqlcode}
\end{frame}
% ---------------------------------------------------------------------
\begin{frame}{Una lista: ¿qué empleados tienen citas?}
\small \texttt{IN} ya lo conoces: \texttt{WHERE id IN (1, 2)}. Ahora la lista la escribe \textbf{una consulta}.
\vspace{0.2cm}
\begin{center}
\begin{tikzpicture}
  \node (ci) at (0,0) {\mini{cita}{rr}{\cab{id} & \cab{empleado\_id}}{%
    1 & 1 \\ 2 & 2 \\ 3 & 1 \\ 4 & 2 \\ 5 & 1 \\}};
  \node (li) at (3.6,0) {\forma{|c|}{\cellcolor{yellow!45}1 \\ \hline \cellcolor{yellow!45}2 \\ \hline}};
  \node[font=\scriptsize\bfseries,above=2pt of li] {la lista};
  \node (em) at (7.6,0) {\mini{empleado}{rll}{\cab{id} & \cab{nombre} & }{%
    \fila 1 & Ana Torres & \si \\
    \fila 2 & Bruno Díaz & \si \\
    3 & Carla Núñez & \no \\}};
  \draw[flecha] (ci.east)--node[above,font=\scriptsize]{\npieza{1}}(li.west);
  \draw[flecha] (li.east)--node[above,font=\scriptsize]{\npieza{2}}node[below,font=\scriptsize]{¿está?}(em.west);
\end{tikzpicture}
\end{center}
\vspace{0.1cm}
\begin{sqlcode}[.62\textwidth]
\kw{SELECT} nombre \kw{FROM} empleado\\
\kw{WHERE} id \kw{IN} \dentro{(\kw{SELECT} empleado\_id \kw{FROM} cita)};
\end{sqlcode}
\vspace{-0.1cm}
\centering\small Resultado: \textbf{Ana Torres} y \textbf{Bruno Díaz}.
\end{frame}
% ---------------------------------------------------------------------
\begin{frame}{\texttt{NOT IN}: lo que no está en la lista}
\small La misma lista, pero ahora buscas a quien \textbf{no} aparece en ella.
\vspace{0.2cm}
\begin{sqlcode}[.62\textwidth]
\kw{SELECT} nombre \kw{FROM} empleado\\
\kw{WHERE} id \kw{NOT IN} \dentro{(\kw{SELECT} empleado\_id \kw{FROM} cita)};
\end{sqlcode}
\vspace{0.1cm}
\begin{columns}[c,onlytextwidth]
\column{.48\textwidth}
\centering
\mini{empleado}{rll}{\cab{id} & \cab{nombre} & }{%
    1 & Ana Torres & \no \\
    2 & Bruno Díaz & \no \\
    \fila 3 & Carla Núñez & \si \\}
\column{.48\textwidth}
\res{l}{\cab{nombre}}{Carla Núñez \\}
\\[8pt]\centering\scriptsize la única que todavía \textbf{no tiene citas}
\end{columns}
\end{frame}
% ---------------------------------------------------------------------
\begin{frame}{Ojo: el signo \texttt{=} espera un solo valor}
\begin{sqlcode}[.8\textwidth]
\kw{SELECT} nombre \kw{FROM} empleado\\
\kw{WHERE} id \mal{=} \dentro{(\kw{SELECT} empleado\_id \kw{FROM} cita)};\\[4pt]
\color{utpRojo}ERROR 1242 (21000): Subquery returns more than 1 row
\end{sqlcode}
\vspace{0.2cm}
\begin{center}
\begin{tikzpicture}[font=\small]
  \node (l) {\forma{|c|}{1 \\ \hline 2 \\ \hline 1 \\ \hline}};
  \node[right=0.6cm of l,font=\Large\ttfamily] (q) {= ?};
  \node[right=0.6cm of q,text width=6cm] {¿Con cuál de todos compara? El gestor no puede elegir y se detiene.};
\end{tikzpicture}
\end{center}
\vspace{0.2cm}
\concepto{Regla}{\small Si la de adentro devuelve \textbf{un valor}, usa \texttt{=}, \texttt{>}, \texttt{<}. Si devuelve \textbf{una lista}, usa \texttt{IN}.}
\end{frame}
% ---------------------------------------------------------------------
\begin{frame}{Una tabla: ¿cuántas citas tiene el que más tiene?}
\small \texttt{MAX(COUNT(*))} \textbf{no se puede} escribir. Primero cuentas; luego buscas el máximo.
\vspace{0.2cm}
\begin{center}
\begin{tikzpicture}
  \node (ci) at (0,0) {\mini{cita}{rr}{\cab{id} & \cab{empleado\_id}}{%
    1 & 1 \\ 2 & 2 \\ 3 & 1 \\ 4 & 2 \\ 5 & 1 \\}};
  \node[draw=utpRojo,very thick,dashed,rounded corners=3pt,fill=yellow!20,inner sep=6pt] (co) at (4.6,0)
    {\mini{conteo}{rr}{\cab{empleado\_id} & \cab{total}}{1 & 3 \\ 2 & 2 \\}};
  \node[font=\scriptsize,text=utpRojo,below=2pt of co] {tabla temporal};
  \node[caja,cajaverde,font=\large\bfseries,minimum width=1.2cm] (r) at (8.6,0) {3};
  \draw[flecha] (ci.east)--node[above,font=\scriptsize]{\npieza{1} contar}(co.west);
  \draw[flecha] (co.east)--node[above,font=\scriptsize]{\npieza{2} \texttt{MAX}}(r.west);
\end{tikzpicture}
\end{center}
\vspace{0.1cm}
\begin{sqlcode}[.62\textwidth]
\kw{SELECT MAX}(total)\\
\kw{FROM} \dentro{(\kw{SELECT} empleado\_id, \kw{COUNT}(*) \kw{AS} total}\\
\ind\ind\dentro{\kw{FROM} cita \kw{GROUP BY} empleado\_id)} \kw{AS} conteo;
\end{sqlcode}
\vspace{-0.3cm}
\ojo{En \texttt{FROM}, la subconsulta \textbf{necesita un nombre} (\texttt{AS conteo}). Sin él: \texttt{ERROR 1248: Every derived table must have its own alias}.}
\end{frame}
% ---------------------------------------------------------------------
\begin{frame}{¿Dónde puede ir una subconsulta?}
\begin{center}
\renewcommand{\arraystretch}{1.5}
\begin{tabular}{@{}>{\bfseries}l>{\raggedright}p{3.4cm}>{\raggedright\arraybackslash}p{7.0cm}@{}}
\toprule
Dónde & \bfseries Qué devuelve & \bfseries Ejemplo \\
\midrule
WHERE & un valor & \footnotesize\ttfamily \kw{WHERE} precio > \dentro{(\kw{SELECT AVG}(precio) \ldots)} \\
WHERE & una lista & \footnotesize\ttfamily \kw{WHERE} id \kw{IN} \dentro{(\kw{SELECT} empleado\_id \ldots)} \\
FROM & una tabla & \footnotesize\ttfamily \kw{FROM} \dentro{(\kw{SELECT} \ldots)} \kw{AS} conteo \\
SELECT & un valor & \footnotesize\ttfamily \kw{SELECT} nombre, \dentro{(\kw{SELECT AVG}(precio) \ldots)} \\
\bottomrule
\end{tabular}
\end{center}
\vspace{0.3cm}
\small En \texttt{SELECT}, el valor aparece como \textbf{una columna más}, repetido en cada fila:
\begin{sqlcode}[.8\textwidth]
\kw{SELECT} nombre, precio, \dentro{(\kw{SELECT AVG}(precio) \kw{FROM} servicio)} \kw{AS} promedio \kw{FROM} servicio;
\end{sqlcode}
\end{frame}
% ---------------------------------------------------------------------
\begin{frame}{¿Subconsulta o \texttt{JOIN}?}
\small Muchas preguntas se pueden responder de las dos formas:
\vspace{0.1cm}
\begin{columns}[T,onlytextwidth]
\column{.49\textwidth}
\centering\etiqueta{Con subconsulta}
\begin{sqlcode}[\dimexpr\textwidth-2\fboxsep\relax]
\kw{SELECT} nombre \kw{FROM} empleado\\
\kw{WHERE} id \kw{IN} (\kw{SELECT} empleado\_id\\
\ind\ind\ind\ind\kw{FROM} cita);
\end{sqlcode}
\column{.49\textwidth}
\centering\etiqueta{Con JOIN}
\begin{sqlcode}[\dimexpr\textwidth-2\fboxsep\relax]
\kw{SELECT DISTINCT} e.nombre\\
\kw{FROM} empleado e\\
\kw{JOIN} cita c \kw{ON} c.empleado\_id = e.id;
\end{sqlcode}
\end{columns}
\vspace{0.1cm}
\centering\small Las dos dan \textbf{Ana Torres} y \textbf{Bruno Díaz}. Sin \texttt{DISTINCT}, el \texttt{JOIN} repite a Ana tres veces.
\vspace{0.3cm}
\raggedright
\begin{itemize}\setlength\itemsep{6pt}
  \item La \textbf{subconsulta} se lee casi como la pregunta: «empleados cuyo id está en las citas».
  \item El \textbf{JOIN} conviene cuando quieres ver columnas de \textbf{las dos} tablas.
\end{itemize}
\end{frame}
% ---------------------------------------------------------------------
\begin{frame}{Receta para escribir una subconsulta}
\begin{center}
\begin{tikzpicture}[font=\small]
  \tikzset{p/.style={caja,text width=3.8cm,minimum height=2.1cm}}
  \node[p] (a) at (0,0) {\npieza{1}\\[3pt]Escribe la consulta \textbf{de adentro} y ejecútala \textbf{sola}.};
  \node[p] (b) at (4.8,0) {\npieza{2}\\[3pt]Mira qué devuelve:\\ ¿un \textbf{valor}, una\\ \textbf{lista} o una \textbf{tabla}?};
  \node[p,cajanaranja] (c) at (9.6,0) {\npieza{3}\\[3pt]Ponla \textbf{entre paréntesis} donde iría ese resultado.};
  \draw[flecha] (a)--(b); \draw[flecha] (b)--(c);
\end{tikzpicture}
\end{center}
\vspace{0.4cm}
\analogia{Si la de adentro funciona sola y devuelve lo que esperas, la consulta completa casi siempre funciona.}
\end{frame}
% =======================================================================
\begin{frame}{\ins{Práctica} Subconsultas en la peluquería}
\small Usa \texttt{ejemplos/subconsultas/3\_practica.sql}. Para cada una, sigue la receta.
\vspace{0.2cm}
\begin{enumerate}\setlength\itemsep{7pt}
  \item ¿Qué servicios cuestan \textbf{menos} que el promedio?
  \item ¿Cuál es el servicio \textbf{más barato}?
  \item ¿Qué mascotas son de \textbf{Rosa Medina}?
  \item ¿Qué clientes tienen al menos un \textbf{perro}?
  \item \textbf{Reto:} ¿qué citas todavía \textbf{no tienen pago}?
\end{enumerate}
\vspace{0.2cm}
\pregunta{\small En cada una: ¿la de adentro devuelve un valor o una lista? ¿Usaste \texttt{=} o \texttt{IN}?}
\end{frame}
\begin{frame}[plain]
  \begin{tikzpicture}[remember picture,overlay]
    \node[anchor=north west,font=\bfseries\fontsize{24}{28}\selectfont]
      at ([xshift=12mm,yshift=-9mm]current page.north west) {Preguntas};
    \node[anchor=north east,inner sep=0pt]
      at ([xshift=-10mm,yshift=-8mm]current page.north east)
      {\includegraphics[width=43mm]{utp-logo.png}};
    \node[anchor=center,inner sep=0pt]
      at ([yshift=-5mm]current page.center)
      {\includegraphics[width=78mm]{preguntas-ilustracion.png}};
  \end{tikzpicture}
\end{frame}
\begin{frame}{Cierre}
\begin{center}
\renewcommand{\arraystretch}{1.35}
\begin{tabular}{>{\raggedright\bfseries}p{3.4cm}>{\raggedright\arraybackslash}p{8.4cm}}
\toprule
Servicio & Programa sin ventana que arranca con Windows: \texttt{mysqld.exe}, \texttt{sqlservr.exe}. \\
Cliente & Cualquier programa que se conecta al gestor: SSMS, Workbench, la terminal, una app. \\
Conectarse & Host, puerto, usuario, contraseña y base. \\
Subconsulta & Una consulta dentro de otra, entre paréntesis. Se resuelve \textbf{primero}. \\
Un valor & Se compara con \texttt{=}, \texttt{>}, \texttt{<}. \\
Una lista & Se usa con \texttt{IN} o \texttt{NOT IN}. \\
Una tabla & Va en \texttt{FROM} y lleva nombre (\texttt{AS}). \\
\bottomrule
\end{tabular}
\end{center}
\end{frame}

\begin{frame}{Referencias}
\footnotesize
\begin{itemize}\setlength\itemsep{8pt}
\item Silberschatz, A., Korth, H. F., \& Sudarshan, S. (2014). \textit{Fundamentos de bases de datos} (6.ª ed.). McGraw-Hill.
\item Oracle. MySQL Reference Manual: \href{https://dev.mysql.com/doc/refman/8.0/en/connecting.html}{\textit{Connecting to the MySQL Server}} y \href{https://dev.mysql.com/doc/refman/8.0/en/subqueries.html}{\textit{Subqueries}}.
\item Microsoft Learn. \href{https://learn.microsoft.com/en-us/sql/ssms/quickstarts/ssms-connect-query-sql-server}{\textit{Connect to and query a SQL Server instance (SSMS)}}.
\end{itemize}
\end{frame}
% Cierre UTP
{
  \setbeamertemplate{background canvas}{%
    \includegraphics[width=\paperwidth,height=\paperheight]{gracias-final.pdf}%
  }
  \setbeamertemplate{footline}{}
  \begin{frame}[plain]
  \end{frame}
}

\end{document}
